\chapter{Symmetry in a lifted template}\label{ch:3}

This chapter holds the theory of the thesis. Its first half explains how a symmetry gets into a lifted template and why it then holds exactly. A template cannot tell apart constants that it does not name (Sections~\ref{sec:3.1} and \ref{sec:3.2}). What the template names and which background facts it is given decide, together, which group the model has (Section~\ref{sec:3.3}). The same theorem shows two different ways of taking a symmetry out again (Section~\ref{sec:3.4}). Section~\ref{sec:3.5} sorts symmetries by what they do to the puzzles a generator produces. The second half puts a price on a symmetry before any training: the data it saves and the ceiling it sets (Sections~\ref{sec:3.6} and \ref{sec:3.7}).

None of the proofs is long. What the chapter adds is that every statement is tied either to a check that needs no training or to a count that was made before the runs of Chapters~\ref{ch:6} and \ref{ch:7}.

\section{Templates, samples and renamings}\label{sec:3.1}

We use the setting of Section~\ref{sec:2.3}. A template T is a set of weighted rules. Most arguments of its atoms are variables, but a rule may also write a constant. The rule body \texttt{col(P, "c1")}, for example, holds only for panels of the first column. We call the constants written in T its \emph{named} constants.

A sample is a set of facts F. Here a sample is one puzzle, or one component of one puzzle (Section~\ref{sec:4.4}), and its facts are of two sorts. The \emph{puzzle facts} describe what the panels show: which objects a panel holds, which component an object belongs to, and the value of each attribute. The \emph{background facts} describe the frame the puzzle sits in, which is the same for every puzzle of a dataset: which panels form a row, which rows are complete, which column a panel stands in, and, for the value-axis model of Section~\ref{sec:4.5}, how the values relate to each other. Grounding T on F gives a network in which every atom a has a value. We write it \(\mathrm{val}_F(a)\).

A \emph{renaming} σ is a one-to-one map from constants to constants. It acts on a fact or an atom by renaming every constant in it, and on a set of facts one fact at a time. A transformation g of a puzzle \emph{is a renaming} if some σ turns the facts of the puzzle into exactly the facts of the transformed puzzle. Exchanging rows 1 and 2 is a renaming: it swaps the panels of the two rows and the row names r1 and r2. Mirroring a left-right puzzle renames the components k0 and k1 into each other. Exchanging size and colour renames the two attributes. Not every transformation is a renaming. Moving an object to another slot of a grid changes a vector of numbers inside a fact, not a constant, and no σ does that.

As in Section~\ref{sec:2.4}, a model has a symmetry \emph{built in} if f(g·x) = g·f(x) for every setting of its weights. For a renaming this has a plain meaning: the output on the transformed puzzle is the output on the original, with its atoms renamed. If no output atom contains a renamed constant, as when rows are exchanged, the symmetry is an invariance. If some do, as the output blocks for size and colour do, it is an equivariance.

\section{The renaming theorem}\label{sec:3.2}

\textbf{Theorem.} Let σ be a renaming that leaves every named constant of T in place. Then for every atom a and for every setting of the weights,

\[ \mathrm{val}_{\sigma(F)}(\sigma(a)) = \mathrm{val}_F(a). \]

\emph{Proof.} A ground rule is a rule r of T whose variables have been replaced by constants through a substitution θ. Since σ leaves the constants of r in place, renaming the ground rule rθ gives r(σ∘θ). That is a ground rule of the same rule r, and its body holds in σ(F) exactly when the body of rθ holds in F. The inverse of σ also leaves the named constants in place, so this pairs the ground rules of the two networks one to one. Weights belong to the rules of T and not to their groundings, so paired ground rules carry the same weights. An aggregation, such as an average or a sum, sees the multiset of values of the bodies that share a head, and σ maps this multiset one to one onto its counterpart. A product or a concatenation combines the atoms of a single body, and σ maps those one by one. For the facts, the equation holds by the definition of σ(F). The ground network has no cycles, so induction over the depth of an atom carries the equation from the facts to every atom. ∎

A transformation of the puzzle is therefore built in whenever two conditions hold:

\begin{enumerate}
\item it is a renaming: some σ turns the facts of the puzzle, background facts included, into exactly the facts of the transformed puzzle;
\item the template names none of the constants that σ moves.
\end{enumerate}

Three remarks.

\begin{itemize}
\item The conditions are sufficient, not necessary. A template can also be symmetric by coincidence, for instance when two rules that name different constants happen to have equal weights.
\item A named constant allows the symmetry to break, but does not force it. The checks below show that with random weights it does break.
\item The theorem is the form that the permutation equivariance of graph networks takes in a logic program. In a graph network the renamed constants are the nodes, and the only facts are edges and node features, which a permutation of the nodes maps onto themselves. In a template any constant can take that role: a row, a component, an attribute, a value. The background facts decide which renamings keep the frame of the puzzle.
\end{itemize}

\textbf{The checks.} Every built-in symmetry claimed in this thesis was also checked numerically, without training. With random weights, several draws and unseen puzzles, the output on the transformed puzzle is compared with the renamed output on the original (\texttt{core/pynl/facet\_symmetry\_check.py}, and the \texttt{--check-shift} and \texttt{--check-scale} options of \texttt{core/pynl/value\_axis.py}). Where the theorem applies, the largest difference lies between 0 and 10⁻¹⁵, the size of a change in the order of a floating-point sum. Where the template names a constant that the transformation moves, the output changes by 0.4 to 1.0. Sections~\ref{sec:3.3} and \ref{sec:3.4} give the individual numbers.

\section{What is named, and what the background facts keep}\label{sec:3.3}

The two conditions of the theorem are two places to put a symmetry. The rules decide which constants are named, and the background facts decide which renamings map a sample onto itself. Together they fix the group of the model: the renamings that leave every named constant in place and keep every background fact true. Adding a background relation can only make this group smaller, because each new fact is one more thing a renaming must keep. The background facts are therefore where a lifted template chooses its group, in the way that the choice of domain fixes the group of a convolutional or a graph network (Section~\ref{sec:2.4}).

The relations used in this thesis show how this works.

\begin{itemize}
\item \textbf{Complete rows}, \texttt{full(r1)} and \texttt{full(r2)}. Both complete rows carry the same fact, and the rules reach them only through it, so exchanging rows 1 and 2 keeps the frame. The question row is set apart, because the rules name it, as in \texttt{inrow(P, "r3")}, and no renaming may move it.
\item \textbf{Columns}, \texttt{col(P, "c1")}, \texttt{col(P, "c2")} and \texttt{col(P, "c3")}. The rules name each column, so no exchange of columns is built in. That is intended: a Progression row read backwards is a different row. For I-RAVEN-X with rows of ten, \emph{column roles} give each panel one of the facts \texttt{col(P, "first")}, \texttt{col(P, "mid")} and \texttt{col(P, "last")} instead. All eight middle columns now carry the same fact, so any exchange of middle columns is built in, while the first and last columns stay apart.
\item \textbf{Neighbours}, \texttt{next(P, Q)}. The row fold of Appendix~\ref{app:E} reads a row as a chain of neighbouring panels. Reversing a row would map \texttt{next} onto its inverse, which is not among the facts, so the fold keeps the direction of a row.
\item \textbf{Offsets on the value axis}, \texttt{off(U, V, k)} for V = U + k. In the value-axis model every value is a constant, and no rule names one. Moving all values by the same amount s maps \texttt{off} onto itself, exactly on a ring of values and away from the ends on a line, so the model is equivariant to this common shift.
\item \textbf{Addition}, \texttt{plus(U, W, V)} for U + W = V. The common shift does not keep these facts, since (u + s) + (w + s) is not (u + w) + s. Giving a model \texttt{plus} therefore takes the common shift out of its group. What remains are the moves that addition keeps, such as adding s to one operand and to the result.
\end{itemize}

\begin{center}\small
\begin{tabularx}{\linewidth}{@{}YYYY@{}}
\toprule
background facts & the renaming they keep & built in & random weights, largest change \\
\midrule
\texttt{full(Rw)} & rows 1 and 2 exchanged & yes & 3·10⁻¹⁶ \\
\texttt{col(P, c)}, one fact per column & no exchange of columns & no & columns 1 and 2: 0.46 to 0.96 \\
\texttt{col(P, "mid")}, column roles & any two middle columns exchanged & yes & 4.4·10⁻¹⁶ (named columns: 0.38) \\
\texttt{off(U, V, k)} & all values moved by the same amount & yes & 3·10⁻¹⁷ \\
\texttt{off} and \texttt{plus} & one operand and the result moved together & the common shift no longer & 0 in 94–97\% of puzzles; up to 2.8·10⁻³ where the shift changes a sum \\
\bottomrule
\end{tabularx}
\end{center}

The same mechanism lets a symmetry be stated as a relation, not only left unnamed. The two components of the left-right and up-down layouts are exchangeable because no rule names k0 or k1. One could also give the model a relation that says how the components of a layout correspond, for example a reflection with one fact per layout, and let every rule take the layout as a variable. The left-right and up-down layouts would then share every weight because no rule names either of them, and the relation would tell the model which component of one layout corresponds to which of the other. The experiments did not need this. It is the form in which the method reaches from the symmetries inside one puzzle to symmetries between layouts.

\section{Two ways to break a symmetry}\label{sec:3.4}

To measure what a symmetry is worth, Chapter~\ref{ch:6} takes it out of a model and compares. The theorem has two conditions, and there is one way to break each.

\textbf{Labels} break the first condition. The template stays as it is, but facts tie each of the exchangeable constants to a fixed vector: \texttt{cid(k0) = [1, 0]} and \texttt{cid(k1) = [0, 1]} for the components, or a one-hot \texttt{aid(A)} for the attributes, read by the summary of every panel. When a puzzle is mirrored, its objects move from k0 to k1, but \texttt{cid(k0)} stays [1, 0]. The facts of the mirrored puzzle are then no longer a renaming of the original facts. With random weights the mirror changes the output by 0.37, and with attribute labels the exchange of type and colour changes it by 0.41. Every weight is still shared between the components or the attributes. The labels add 320 weights to the 266,528 of the components model, or 512 to the 288,328 of the attributes model.

\textbf{Separate weights} break the second condition. Every rule with weights is written once per component, or once per attribute, with that constant in it. The template now names the constants, and the pairing of ground rules in the proof is gone: the rules for k0 and those for k1 have weights of their own. With random weights the mirror changes the output by 0.78, and the exchange of type and colour by 0.90. At the same widths the model has twice as many weights for the components and 3.2 times as many for the attributes, so Chapter~\ref{ch:6} also uses narrower versions with as many weights as the symmetric model.

The two ways take out different things. Labels remove the exact invariance and keep the sharing: every training puzzle still trains every weight. Separate weights remove both. Comparing a symmetric model with its labelled and its separate versions therefore tells apart what the invariance is worth and what the sharing is worth (Section~\ref{sec:6.3}).

A constant can also be named by accident. The first version of the copy links in Section~\ref{sec:6.5} gave every panel a gate of its own, \texttt{gate("p3")}. That names the panels, so the exchange of rows 1 and 2 was no longer built in (0.99 with random weights), and the trained model made use of the freedom. Gates computed from what a panel contains name nothing, and with them the exchange is exact again (4.4·10⁻¹⁶).

\section{Three kinds of symmetry}\label{sec:3.5}

Whether a built-in symmetry helps or hurts depends on what it does to the puzzles that the generator produces. Take a transformation g of a complete matrix that leaves the missing panel in its place, a valid puzzle x and its answer y. There are three cases.

\begin{itemize}
\item \textbf{A data symmetry.} g·x is again a valid puzzle, its answer is g·y, and the generator produces g·x as often as x. The training data then contains both, and a model with the symmetry has less to learn. This is a case of what \citet{wang2023general} call correct equivariance; their definition asks only that g·x occur in the data, not that it occur as often. Section~\ref{sec:3.6} counts what it is worth.
\item \textbf{Out of support.} The generator never produces g·x, so the symmetry puts no constraint on the data. This is Wang et al.'s extrinsic equivariance. For the benchmark two cases of it differ. If g·x is still a valid puzzle with the answer g·y, the symmetry belongs to the task but not to the data. Building it in costs nothing on the right solution, and it rules out solutions that hold only on what the generator produces (Section~\ref{sec:6.5}). If g·x is not a valid puzzle at all, the model never has to answer it, so no ceiling follows, but the model must find the rule in a form that respects g, and that can cost data (Section~\ref{sec:6.7}).
\item \textbf{Conflicting.} g·x is a valid puzzle that the generator does produce, but its answer is not g·y. A model with the symmetry has to answer x and g·x consistently, so it is wrong on one of them. This is Wang et al.'s incorrect equivariance, for which they prove a lower bound on the error, orbit by orbit. Section~\ref{sec:3.6} evaluates that bound on a rule's table, and Sections~\ref{sec:6.7} and \ref{sec:7.4} test it.
\end{itemize}

Wang et al. make their definitions point by point, so one move can be correct on some rules and incorrect on others, as the shift is.

Which case a transformation falls into is a property of the generator. It can be read off the generators' code and checked on the data. The table gives the share of rule instances (one attribute of one puzzle) that still satisfy a rule of the same kind after the transformation, in the order Constant / Progression / Arithmetic / Distribute. Before the transformation every instance satisfies its own rule.

\begin{center}\small
\begin{tabularx}{\linewidth}{@{}YYYY@{}}
\toprule
transformation & I-RAVEN & I-RAVEN-X, rows of 3 & I-RAVEN-X, rows of 10 \\
\midrule
rows 1 and 2 exchanged & 100 / 100 / 100 / 100 & 100 / 100 / 100 / 100 & 100 / 100 / 100 / \textbf{0} \\
columns 1 and 2 exchanged & 100 / \textbf{0} / 51 / 100 & 100 / \textbf{0} / 55 / 100 & 100 / \textbf{0} / 50 / \textbf{0} \\
two middle columns exchanged & – & – & 100 / \textbf{0} / 100 / \textbf{0} \\
values reflected, v → max − v & 100 / 100 / \textbf{0} / 100 & 100 / 100 / \textbf{0} / 100 & 100 / 100 / \textbf{0} / 100 \\
matrix transposed & 58 / 3 / 1 / 100 & – & – \\
\bottomrule
\end{tabularx}
\end{center}

The symmetries that this thesis builds in, breaks or adds on purpose fall into the three cases as follows.

\begin{center}\small
\begin{tabularx}{\linewidth}{@{}YYYY@{}}
\toprule
symmetry & case & the reason, from the generator & used in \\
\midrule
order of the objects in a panel; order of the candidates & data & the order carries no meaning & every model \\
the two components of the left-right and up-down layouts & data & identical constraints, and matching rule counts and means in both components & 6.2, 6.3 \\
size and colour exchanged (I-RAVEN-X) & data & both are drawn by the same function, with the same rules and range & 6.2, 6.3 \\
the attributes exchanged, apart from size and colour in I-RAVEN-X & out of support & the attributes differ in their values (in I-RAVEN they have 5, 6 and 10) and type never carries Arithmetic & 6.4 \\
rows 1 and 2 exchanged, I-RAVEN and rows of 3 & data & each row is drawn afresh by the same rules, with one small exception for Arithmetic on colour & 6.2 (the vision transformer), 6.5 \\
rows 1 and 2 exchanged, rows of 10 & out of support, valid & a Distribute row is the first row rotated by a step, and after the exchange no step fits; the answer is still the value that the third row lacks & 6.5 \\
middle columns exchanged, rows of 10 (column roles) & out of support: valid for Arithmetic and for the answer of Distribute, not for Progression & the operands of a sum may be reordered and the value a row lacks does not depend on the order, but the generator places the zeros among the operands in a fixed pattern & 6.6 \\
columns 1 and 2 tied & out of support, not valid & a Progression row read in another order is not a Progression row & 6.7 \\
values reflected (a model that cannot tell v from max − v) & conflicting & Constant, Progression and Distribute rows stay valid and equally likely with the reflected answer, which differs & 6.7 \\
all values moved by the same amount & data for Constant, Progression, Distribute and midpoint, exactly on a ring of values; conflicting for sum and Arithmetic & a shifted sum row is not a sum row, but its question row is a valid question with another answer & 6.2, 7.3–7.6 \\
\bottomrule
\end{tabularx}
\end{center}

\section{Pricing by orbits}\label{sec:3.6}

A symmetry merges inputs into orbits. This section counts what that is worth on a rule's table, in the simplest setting where the count can be made exactly: a single row of three values. First the words.

A \textbf{row} is three values (a, b, c). In the \textbf{question row} the third is missing, and a \textbf{learner} is anything that turns the two values (a, b) into an answer for c. We write f(a, b) for that answer.

A \textbf{move} changes a row in a fixed way. Three moves matter here:

\begin{itemize}
\item the \textbf{shift}: add the same number s to all three values;
\item \textbf{scaling}: multiply all three values by the same number k;
\item \textbf{one operand with the answer}: add s to a and to c, or to b and to c.
\end{itemize}

A rule \textbf{keeps} a move if every row of the rule is again a row of the rule after the move. A learner \textbf{has} a move if moving its input moves its answer in the same way: for the shift, f(a + s, b + s) = f(a, b) + s. The \textbf{orbit} of an input is everything the moves turn it into. Values are the whole numbers 0 to 50. Where a statement divides, it is made on the ring of 101 values that the models use, where every step has a half and every multiple.

\subsection{Which rule keeps which move}\label{sec:3.6.1}

\textbf{Statement 1.} A rule of the form c = pa + qb keeps the shift exactly when p + q = 1. Every such rule keeps scaling. Sum keeps the move of one operand with the answer, and midpoint does not.

\emph{Why.} Shift a row of the rule by s. The new third value is c + s. The rule asks for p(a + s) + q(b + s), which is pa + qb + (p + q)s, that is c + (p + q)s. The two agree for every s exactly when p + q = 1. Midpoint has ½ + ½, mirror (c = 2a − b) has 2 − 1, Progression (c = 2b − a) has −1 + 2 and Constant (c = b) has 0 + 1: all keep the shift. Sum has 1 + 1 = 2 and does not. For scaling, multiply the rule by k: kc = p(ka) + q(kb), which holds whatever p and q are. For the last move, add s to a and to c in a sum row: c + s = (a + s) + b holds; in a midpoint row the rule asks for c + s/2.

\emph{Checked on the tables} (rows that are rows again after the move, of those that stay in range):

\begin{center}\small
\begin{tabularx}{\linewidth}{@{}YllY@{}}
\toprule
 & shift & scaling by 2, 3, 5 & one operand with the answer \\
\midrule
sum, 1,225 rows & 0 of 30,325 & 465 of 465 & 39,200 of 40,425 \\
midpoint, 1,250 rows & 40,400 of 40,400 & 490 of 490 & 0 of 51,450 \\
mirror, 1,250 rows & 40,400 of 40,400 & 490 of 490 & 0 of 51,450 \\
\bottomrule
\end{tabularx}
\end{center}

The 3\% that sum loses under its own move are rows pushed onto a = 0, which the table leaves out.

\subsection{A learner with the shift, and what it saves}\label{sec:3.6.2}

\textbf{Statement 2.} A learner that has the shift can only compute c = b + g(b − a), for some list of numbers g with one entry per first step b − a.

\emph{Why.} Shift the input by −a. The learner must answer (0, b − a) with f(a, b) − a. So f(a, b) = a + f(0, b − a): its answer is a plus something that depends on the first step alone. Writing g(d) for f(0, d) − d gives f(a, b) = b + g(b − a). All inputs with the same first step form one orbit, and the learner has one free number for each orbit.

\textbf{Statement 3 (the data saved).} If the rule keeps the shift, one training row teaches the learner every row with the same first step. With rows drawn evenly from a table in which orbit o holds a share P(o) of the rows, the expected share of unseen rows whose orbit n training rows have reached is

the sum over orbits of P(o) · (1 − (1 − P(o))ⁿ).

\emph{Why.} A rule that keeps the shift is itself of the form c = b + g*(b − a), by the argument of Statement 2 applied to the rule. A training row with first step d shows g*(d), and that is the right entry for every row with that first step. An unseen row is therefore answered correctly exactly when its first step occurred in training. A given orbit is missed by one training row with chance 1 − P(o) and by n independent rows with chance (1 − P(o))ⁿ. An unseen row lies in orbit o with chance P(o). Adding up over the orbits gives the formula. The same holds for any set of moves the rule keeps, with its own orbits.

\emph{Checked on midpoint.} ``Expected'' uses the table alone. ``Reached'' is the counting learner on the actual training puzzles, which keeps one answer per orbit and is told the rule. 100 and 600 training puzzles hold 57 and 311 midpoint question rows.

\begin{center}\small
\begin{tabularx}{\linewidth}{@{}YlYYl@{}}
\toprule
built in & orbits & 100 puzzles: expected / reached / trained network & 600 puzzles & rows for 90\% \\
\midrule
nothing, one answer per pair (a, b) & 1,250 & 4.5 / 5.0 / 5.0 ± 0.5 & 22.0 / 21.6 / 18.0 ± 0.9 & 2,878 \\
the shift & 50 & 74.8 / 80.0 / 59.9 ± 4.5 & 98.7 / 99.8 / 98.3 ± 1.9 & 107 \\
the shift and doubling & 26 & 90.5 / 94.3 / 89.2 ± 2.7 & 99.3 / 99.8 / 99.2 ± 0.2 & 55 \\
the shift and scaling & 1 & 100 / 100 / 100.0 & 100 / 100 / 100.0 & 1 \\
\bottomrule
\end{tabularx}
\end{center}

The formula gives the counting learner's score to within 6 points before a single puzzle is drawn. The shift cuts the rows needed for 90\% from 2,878 to 107, a factor of 27. The trained networks are within 4 points of the counted value without a symmetry. With one they are under it where many orbits are left and little data is given: by 20 points with the shift and by 5 with doubling at 100 puzzles, by 2 or less at 600, and by nothing with scaling. The three networks with a symmetry are those whose row descriptor reads the steps (Section~\ref{sec:7.7}); Section~\ref{sec:3.7} says why that matters.

\subsection{The ceiling, and whose ceiling it is}\label{sec:3.6.3}

\textbf{Statement 4 (the ceiling).} On a rule in which the rows of an orbit all need different answers, a learner of the form of Statement 2 is right on at most one row of each orbit. For sum under the shift that is 97 rows of 1,225: 7.9\%.

\emph{Why.} On a sum row the answer is b + a, so the learner is right exactly when g(b − a) = a. Rows with the same first step have different a, and g holds one number for that step. It can match one of them. Sum's table has 97 different first steps. With rows drawn evenly, the best a learner can expect is one right row per orbit, 97 in 1,225. In general the count per orbit is the largest group of its rows whose answers agree once the move is undone; for the shift, the rows with the same c − b. Written this way a learner with the shift is a learner of c − b that ignores the shift, and the count is the bound of \citet[Theorem 4.3]{wang2023general} for such learners, evaluated here on a table.

\emph{Checked.} Fitted on the whole table, the counting learner with the shift is right on 97 of 1,225 rows. With the shift and doubling it is 88 (7.2\%), with the shift and scaling 17 (1.4\%). The trained shared model reaches 4.5 ± 2.3 and 6.2 ± 0.8 at 100 and 600 puzzles and 7.1 at 6,000, at the epoch picked on the selection puzzles.

\emph{How sharp the check is.} The ceiling is on the expected accuracy. The unseen set has 507 sum rows, on which one standard deviation of a 7.9\% rate is 1.2 points. Over all 172 evaluated epochs of the nine runs, 8 are above 7.9 and the highest is 9.7, in the run on 6,000 puzzles, whose epochs average 7.9. So the model sits at the ceiling, within what chance allows on this many rows, and the measurement cannot exclude a point or two above it. Statement 5 says why a point or two would not contradict anything.

\textbf{Statement 5 (the ceiling belongs to learners that answer from the question row).} A model that has the shift of the whole puzzle, all nine panels together, is not capped on sum by that symmetry alone.

\emph{Why.} Here is a function of the whole puzzle that uses only differences:

c₃ = b₃ + (c₁ − b₁) + (a₃ − a₁).

Shift every panel by s and the answer moves by s, so the function has the shift. On a sum puzzle c₁ − b₁ is a₁, so the function returns b₃ + a₁ + a₃ − a₁ = a₃ + b₃, the right answer. A complete sum row shows where zero is, and a learner that can measure from a panel of one row to a panel of another can use it.

\emph{Checked.} The function is right on 507 of 507 unseen sum rows and follows a shift of the whole puzzle on all of them.

\emph{What this means for the trained model.} Statements 2 to 4 are proved for a learner that sees the question row. The template sees the whole puzzle. It stays under the ceiling because of how it reads the puzzle, which can be read off its rules: the score of a candidate value V is computed from V's offsets to the two panels of the question row, from how many panels of each row sit exactly at V, and from a descriptor of the puzzle that is the same for every V. No rule measures the offset between a panel of the question row and a panel of another row. So apart from values that happen to coincide between rows, the template is a learner of the form of Statement 2, with a g that may depend on the descriptor. The rows of an attribute are drawn independently, so the descriptor of the complete rows says nothing about where the question row sits. For the template the ceiling is therefore a prediction that follows from its structure, and not a consequence of the puzzle-wide shift alone. The coincidences are a real way past it, though a small one: in 48 of the 507 unseen sum rows the answer is a value that already stands in the puzzle, and the template sees which values those are. The runs sit at the ceiling and not clearly above it.

In the terms of \citet{wang2023general}: shifting one row turns a sum row into a row that needs another answer, which is their incorrect equivariance, the case their lower bound covers. Shifting a whole sum puzzle turns it into something that is not a sum puzzle at all, which is their extrinsic equivariance. On such orbits only one puzzle occurs in the data and their bound is zero. Extrinsic equivariance can still hurt (\citealp[Section 5]{wang2023general}) or help (\citealp{wang2023surprising}).

\subsection{More of the rule's moves}\label{sec:3.6.4}

\textbf{Statement 6 (doubling).} A learner that has the shift and doubling has one free number for each odd part of the first step. On the ring, a row with a first step it never saw is answerable exactly when a training row had a first step with the same odd part.

\emph{Why.} By Statement 2 the learner computes b + g(d). Doubling the input must double the answer's offset from b, so g(2d) = 2g(d). Write a step as d = m·2ʲ with m odd. Then g(d) = 2ʲ·g(m): the entry for the odd part fixes the entries for all its doublings, and one training row with any of them fixes g(m). Midpoint's even first steps between −50 and 50 have 26 odd parts.

\emph{Checked.} At 100 puzzles 101 unseen midpoint rows have a new first step: 72 with an odd part that training had, 29 without. The trained doubled model is right on 50 / 43 / 39 of the 72 and on 0 / 0 / 0 of the 29, over three seeds. The model with the shift alone is right on none of the 101.

\textbf{Statement 7 (scaling).} On the ring, a learner that has the shift and every scaling computes c = a + t·(b − a) for one number t, on every row with a ≠ b.

\emph{Why.} Any two different values (a, b) are the image of (0, 1) under ``multiply by b − a, then add a''. The learner has both moves, so f(a, b) = (b − a)·f(0, 1) + a. Its only free number is t = f(0, 1). Midpoint is t = ½, Progression is t = 2, mirror is t = −1.

\emph{A consequence that is also proved.} With scaling built in, every question row with two different values is the same input, (0, 1). The question row then says nothing about which rule applies. Midpoint and Progression need different t, so a model with scaling can tell them apart only from the complete rows.

\emph{Checked.} The counting learner with these moves is right on every midpoint row after the first training row. On sum its ceiling is 17 rows of 1,225. The consequence is what Section~\ref{sec:7.7} measured, in both directions. With a row descriptor that came out the same for every row, the scaled model continued four midpoint rows in five as a Progression at two seeds of three. With a descriptor that reads each pair's step, it is right on all 505 midpoint rows in six runs of six, from 100 puzzles and from 600, and the descriptor separates the two rules in every run.

\textbf{Statement 8 (addition's own moves).} A learner that has both moves of one operand with the answer computes c = a + b + k for one number k.

\emph{Why.} Move the first operand by −a and the second by −b: the learner must answer (0, 0) with f(a, b) − a − b. So f(a, b) = a + b + f(0, 0).

\emph{Checked.} The counting learner with these moves is right on every sum row after one training row. The shared model given the \texttt{plus} facts reaches 98.2 ± 0.8 on sum from 100 puzzles and 99.7 ± 0.4 from 600.

\textbf{Statement 9 (one orbit is the rule, up to one answer).} If the moves a learner has can turn any input into any other, the learner is fixed by its answer to one input.

\emph{Why.} Pick an input x. Every other input is some move h applied to x, and the learner must answer it with h applied to its answer to x. Nothing else is free.

This is why the number of orbits is the right unit: it counts the answers that are left to learn. Statements 7 and 8 are the two cases with one orbit, and in both the model has been handed its rule up to one number.

\subsection{Values read as numbers}\label{sec:3.6.5}

\textbf{Statement 10.} Let the rule be c = pa + qb, and let the two complete rows of a puzzle be (a₁, b₁, c₁) and (a₂, b₂, c₂). If a₁b₂ ≠ a₂b₁, the two rows fix p and q, a least-squares fit finds them, and its answer for the question row is exact. Nothing is learned from other puzzles.

\emph{Why.} The two rows are two equations, a₁p + b₁q = c₁ and a₂p + b₂q = c₂. They have exactly one solution when a₁b₂ ≠ a₂b₁, and the rule's own p and q are a solution. Least squares returns the exact solution when there is one. When a₁b₂ = a₂b₁ the rows leave a whole line of solutions, and the fit returns the shortest, which need not be the rule's.

\emph{Checked} (\texttt{analysis/twins\_linear\_fit.py}, unseen puzzles, no training):

\begin{center}\small
\begin{tabular}{@{}lllll@{}}
\toprule
puzzles & the fourth rule & Progression & Constant & Distribute Three \\
\midrule
sum & 99.6 & 98.7 & 100 & 0 \\
midpoint & 99.8 & 99.3 & 100 & 0 \\
mirror & 99.8 & 98.7 & 100 & 0 \\
lookup & 1.7 & 98.6 & 100 & 0 \\
\bottomrule
\end{tabular}
\end{center}

All 44 misses on the linear rules (13, 7, 12 and 12) are puzzles with a₁b₂ = a₂b₁. There is no miss among the others. Lookup and Distribute Three are not of the form, and the fit gets 1.7\% and nothing.

\subsection{Reach: what is proved and what is measured}\label{sec:3.6.6}

\textbf{Statement 11.} For a learner with one free number per orbit, training rows from other orbits put no condition on the number of an orbit that training never reached.

\emph{Why.} By Statement 2 the learner's answer on a row depends on g at that row's first step only. A training row with another first step involves another entry of g. So every setting of the unreached entry fits the training rows equally well, and the training data cannot make it right.

\emph{Proved for:} the counting learner. \emph{Measured for the networks:} the kernel has one free weight vector per offset, and nothing ties the vectors of different offsets together. The template therefore has no way to carry ``half the step back'' from the steps it has seen to one it has not, beyond coincidences of weights trained in other roles; Statement 11 applies to it nearly as it stands. The measurement is a check of that, not a discovery about networks. What the count adds is that it says before training which rows those are. Pooled over three seeds:

\begin{center}\small
\begin{tabularx}{\linewidth}{@{}Ylll@{}}
\toprule
training rows in the row's orbit & 0 & 1 & 3–4 \\
\midrule
shared model, 100 puzzles & 0 of 303 & 422 of 573 (74\%) & 172 of 186 (92\%) \\
shared model, 600 puzzles & 0 of 3 & 71 of 72 (99\%) & 89 of 90 (99\%) \\
weights per value, 100 puzzles (orbit: the pair) & 11 of 1,440 (1\%) & 65 of 75 (87\%) & – \\
\bottomrule
\end{tabularx}
\end{center}

With no training row in its orbit a row is never right, beyond the 1\% that guessing among 101 values gives. With one it usually is. So which unseen rows a model of this family answers can be listed from the training set before training: the rows whose orbit it reached.

\section{The ladder}\label{sec:3.7}

Statements 2, 6, 7 and 10 put the models in one order, by how much of the rule's group is built in and how many answers are left to learn. Midpoint at 100 training puzzles:

\begin{center}\small
\begin{tabularx}{\linewidth}{@{}YYYlY@{}}
\toprule
built in & left to learn & expected from the table & counting learner & trained model \\
\midrule
nothing about the values & 1,250 rows & 4.5 & 5.0 & 5.0 ± 0.5 \\
the shift & 50 orbits & 74.8 & 80.0 & 59.9 ± 4.5 \\
the shift and doubling & 26 orbits & 90.5 & 94.3 & 89.2 ± 2.7 \\
the shift and scaling & 1 orbit, one number & 100 & 100 & 100.0 / 100.0 / 100.0 \\
the linear form, values as numbers & nothing from other puzzles & – & – & 99.8, no training \\
\bottomrule
\end{tabularx}
\end{center}

For sum the ladder has other rungs, because its moves are others: nothing (4.8 ± 0.6); the shift, with the answer read from the question row, which sum does not keep (4.5 ± 2.3 at a ceiling of 7.9); its own moves (98.2 ± 0.8); the linear form (99.6, no training).

The count says each rung needs less data than the one below, and the trained models end in the same order. They do so only at the second attempt, and they do not reach the count's numbers on every rung. With the first row descriptor the two scaling rungs fell below the rung beneath them: 49.9 ± 5.3 with doubling and 19.0 / 100.0 / 21.0 with scaling. Statement 7's consequence names the reason: with rows read in a unit the rule has to come from the complete rows, and that descriptor did not carry it. A descriptor that reads each pair's step was then built, with its outcome predicted before the runs, and gives the table above (Section~\ref{sec:7.7}). The first attempt is a result in its own right: the same group and the same orbit count gave 20\% or 100\% depending on how a row was described. The count is a bound for a learner that shares nothing across orbits and can tell every orbit apart. It predicts the order of the rungs and which rows a model can get, not the accuracy a particular network reaches; with the shift alone the network stays 20 points under the counting learner at 100 puzzles (59.9 against 80.0), and we did not find the cause.

\subsection{What is not proved}\label{sec:3.7.1}

\begin{itemize}
\item \textbf{That the trained network reaches the count.} The count is what a learner with one answer per orbit gets. The networks come within 4 points of it without a symmetry, and with the shift they stay 20 points under it at 100 puzzles and 2 at 600. A network reaches the count only if it can tell the orbits apart, which a row descriptor that pools away the steps prevents (Section~\ref{sec:7.7}). The 20 points with the shift alone may be a limit of the same kind; it is not explained here.
\item \textbf{That a network cannot generalise past its orbits.} Statement 11 is proved for the counting learner. For the networks it is a measurement on one family of templates, in which nothing but the built-in moves ties weights together. A model with other ties, such as a smooth kernel, could do otherwise.
\item \textbf{The ceiling for the template.} Statement 5: it follows from how the template reads the puzzle, with coincidences of values between rows left over. It is not a theorem about the puzzle-wide shift. The runs sit at it: 7.1 or less at every picked epoch, and above 7.9 at 8 of 172 evaluated epochs, 9.7 at the most.
\item \textbf{Anything about groups on 0 to 50.} The models compute on a ring of 101 values, chosen because 101 is prime: every step then has a half and every multiple, so the scaling moves are defined. On the values 0 to 50 the moves are not a group at all, since a shift or a scaling can leave the range. The data never leaves the part where the ring and the whole numbers agree: a midpoint row has a + b even and at most 100, so half of a + b on the ring is the ordinary midpoint. The counts are made on the actual tables. Statements 6 and 7, and the exactness checks of the scaled models (Section~\ref{sec:7.7}), hold on the ring only; on the data, scaling maps most rows out of range and is an out-of-support symmetry in the sense of Section~\ref{sec:3.5}.
\end{itemize}

\section{What is known already, and what is added here}\label{sec:3.8}

The sources were found by search on 2026-11-17; the papers cited for a result were read on 2026-11-23.

\begin{center}\small
\begin{tabularx}{\linewidth}{@{}lYY@{}}
\toprule
statement & standing & source \\
\midrule
1 & textbook: a weighted sum of points needs no origin when the weights sum to 1 & Section~\ref{sec:2.4} (\citealp{stevens1946theory}; \citealp{baez2009torsors}) \\
2, 6 to 9 & elementary; no novelty claimed. The form in 8 is a known functional equation & reading list: \citealp{reem2017remarks} \\
3 & the gain from an invariance as more effective samples is known for kernel methods: for a finite group the effective number of samples is multiplied by the group's size; strict gains and improved bounds of other forms are known in general & \citealp{tahmasebi2023exact}; \citealp{bietti2021sample}; other forms: \citealp{elesedy2021provably}; \citealp{sannai2021improved} \\
3, added here & the exact count on a finite table with orbits of unequal size, and its test against a counting learner and trained networks & – \\
4 & the bound: Theorem 4.3, applied to the answer relative to the moved input & \citealp{wang2023general} \\
5 & the two cases are Wang et al.'s incorrect and extrinsic equivariance; the example and the reading of the template are added here & \citealp{wang2023general} \\
10 & models that make addition an operation of the representation & Section~\ref{sec:2.5} (NVSA, ARLC) \\
11 & what trained networks do on inputs outside the training support is studied for Boolean functions (they find the interpolator of lowest degree); reasoning over symbols never seen in training is proved to fail for fully connected networks & \citealp{abbe2023generalization}; \citealp{boixadsera2024when} \\
\bottomrule
\end{tabularx}
\end{center}
